% Punkte und Strecken im Koordinatensystem – bank/punkte-und-strecken-im-koordinatensystem/ (zusammenbau v0.4, Heft)
\einheitenkopf[t1]{Punkte und Strecken im Koordinatensystem}
\setcounter{aufgabe}{0}

% Anlauf (ohne Prüfkennung)
\begin{aufgabe}{Darstellen und Lage lesen – Anlauf}
\begin{teile}
\teil $P(4 | -1 | 0)$ – was sagt die Null? \\ \kreuz{$P$ liegt auf der $x_3$-Achse} \\ \kreuz{$P$ liegt in der $x_1x_2$-Ebene} \\ \kreuz{$P$ liegt in der $x_2x_3$-Ebene}
\teil $A(5 | 3 | 1)$ und $B(2 | 1 | 4)$ eintragen, $C$ ablesen.

\begin{ksys3} \rpunkt{4}{5}{2}{C} \end{ksys3}
C( \leerfeld | \leerfeld | \leerfeld )
\teil Zwei Masten stehen senkrecht in $F_1(-2 | -2 | 0)$ ($4$ m hoch) und $F_2(5 | 3 | 0)$ ($3$ m hoch); ein Seil verbindet die Spitzen (1 LE = 1 m). Zeichne Masten und Seil.

\begin{ksys3} \end{ksys3}
\end{teile}
\end{aufgabe}

\begin{aufgabe}{Darstellen und Lage lesen – Prüfungsaufgaben}
\begin{teile}
\teil Ein Holzwürfel mit Kante $4$ hat eine Ecke im Ursprung und liegt im ersten Oktanten (alle Koordinaten $\ge 0$). Zeichne das Viereck $P(4 | 0 | 1)$, $Q(1 | 4 | 0)$, $R(0 | 4 | 2)$, $S(3 | 0 | 4)$ ein. \hfill (Abitur 2019 GK)

\begin{ksys3} \rquader{0,0,0}{4}{4}{4} \end{ksys3}
\teil Ein Glaswürfel mit Kante $3$ hat eine Ecke im Ursprung und liegt im ersten Oktanten (alle Koordinaten $\ge 0$). Zeichne das Viereck $P(3 | 1 | 0)$, $Q(3 | 3 | 2)$, $R(1 | 3 | 3)$, $S(0 | 0 | 2)$ ein. \hfill (Abitur 2019 GK)

\begin{ksys3} \rquader{0,0,0}{3}{3}{3} \end{ksys3}
\teil Ein Würfel mit Kante $4$ hat eine Ecke im Ursprung und liegt im ersten Oktanten (alle Koordinaten $\ge 0$). Zeichne das Viereck $P(0 | 3 | 0)$, $Q(4 | 4 | 1)$, $R(4 | 1 | 4)$, $S(1 | 0 | 4)$ ein. \hfill (Abitur 2019 GK)

\begin{ksys3} \rquader{0,0,0}{4}{4}{4} \end{ksys3}
\teil Ein Quader hat eine Ecke im Ursprung $O$ und die Kanten $5$, $6$ und $3$ entlang der $x_1$-, $x_2$- und $x_3$-Achse. $P(5 | 2 | 3)$ liegt auf einer Kante. Zeichne das Dreieck mit den Ecken $O$, $P$ und $R(0 | 6 | 3)$ ein. \hfill (Abitur 2026 GK)

\begin{ksys3} \rquader{0,0,0}{5}{6}{3} \end{ksys3}
\teil Ein Quader hat eine Ecke im Ursprung $O$ und die Kanten $6$, $3$ und $4$ entlang der $x_1$-, $x_2$- und $x_3$-Achse. $P(6 | 0 | 1)$ liegt auf einer Kante. Zeichne das Dreieck mit den Ecken $O$, $P$ und $T(0 | 3 | 4)$ ein. \hfill (Abitur 2026 GK)

\begin{ksys3} \rquader{0,0,0}{6}{3}{4} \end{ksys3}
\teil Ein Quader hat eine Ecke im Ursprung $O$ und die Kanten $3$, $6$ und $3$ entlang der $x_1$-, $x_2$- und $x_3$-Achse. $P(3 | 2 | 3)$ liegt auf einer Kante. Zeichne das Dreieck mit den Ecken $O$, $P$ und $T(3 | 6 | 0)$ ein. \hfill (Abitur 2026 GK)

\begin{ksys3} \rquader{0,0,0}{3}{6}{3} \end{ksys3}
\end{teile}
\end{aufgabe}

\begin{aufgabe}{Darstellen und Lage lesen – Prüfungsaufgaben – weiter}
\begin{teile}
\teil Begründe: $P(-3 | 4 | 3)$ und $Q(5 | -1 | 2)$ liegen auf derselben Seite der $x_1x_2$-Ebene. \hfill (Abitur 2025 GK)
\teil Begründe: Das Dreieck $A(6 | 0 | 1)$, $B(6 | 3 | 5)$, $C(6 | -3 | 5)$ liegt parallel zur $x_2x_3$-Ebene und symmetrisch zur $x_1x_3$-Ebene. \hfill (Abitur 2026 GK)
\teil Begründe: $R(1 | -2 | -3)$ und $T(4 | 5 | 2)$ liegen auf verschiedenen Seiten der $x_1x_2$-Ebene. \hfill (Abitur 2025 GK)
\teil Der Grundriss eines Pavillons ist ein Achteck in der $x_1x_2$-Ebene, symmetrisch zur $x_1$- und zur $x_2$-Achse. Gegeben sind $A(6 | 0)$, $B(5 | 3)$, $C(0 | 4)$. Zeichne das vollständige Achteck. \hfill (Abitur 2025 GK)

\begin{ksys}[xmin=-7,xmax=7,ymin=-7,ymax=7,xlabel=$x_1$,ylabel=$x_2$] \punkt{6}{0}{A} \punkt{5}{3}{B} \punkt{0}{4}{C} \end{ksys}
\teil Der Beet ist ein Achteck in der $x_1x_2$-Ebene, symmetrisch zur $x_1$- und zur $x_2$-Achse. Gegeben sind $A(4 | 0)$, $B(3 | 2)$, $C(0 | 3)$. Zeichne das vollständige Achteck. \hfill (Abitur 2025 GK)

\begin{ksys}[xmin=-7,xmax=7,ymin=-7,ymax=7,xlabel=$x_1$,ylabel=$x_2$] \punkt{4}{0}{A} \punkt{3}{2}{B} \punkt{0}{3}{C} \end{ksys}
\teil Der Bodenplatte ist ein Achteck in der $x_1x_2$-Ebene, symmetrisch zur $x_1$- und zur $x_2$-Achse. Gegeben sind $A(6 | 0)$, $B(2 | 5)$, $C(0 | 6)$. Zeichne das vollständige Achteck. \hfill (Abitur 2025 GK)

\begin{ksys}[xmin=-7,xmax=7,ymin=-7,ymax=7,xlabel=$x_1$,ylabel=$x_2$] \punkt{6}{0}{A} \punkt{2}{5}{B} \punkt{0}{6}{C} \end{ksys}
\end{teile}
\end{aufgabe}

\begin{aufgabe}{Darstellen und Lage lesen – Prüfungsaufgaben – weiter}
\begin{teile}
\teil Eine Pyramide hat den quadratischen Boden $ABCD$ mit der Seite $3$ cm; die Spitze $S$ liegt $4$ cm senkrecht über $A$. Im Gitter sind der Boden und die Spitze des Seitendreiecks $ABS$ markiert. Zeichne die drei übrigen Seitendreiecke ein. \hfill (Abitur 2024 GK)

\begin{ksys}[xmin=-5,xmax=9,ymin=-5,ymax=9] \punkt{0}{0}{A} \punkt{3}{0}{B} \punkt{3}{3}{C} \punkt{0}{3}{D} \punkt{0}{-4}{S} \end{ksys}
\teil Eine Pyramide hat den quadratischen Boden $ABCD$ mit der Seite $4$ cm; die Spitze $S$ liegt $3$ cm senkrecht über $A$. Im Gitter sind der Boden und die Spitze des Seitendreiecks $ABS$ markiert. Zeichne die drei übrigen Seitendreiecke ein. \hfill (Abitur 2024 GK)

\begin{ksys}[xmin=-4,xmax=10,ymin=-4,ymax=10] \punkt{0}{0}{A} \punkt{4}{0}{B} \punkt{4}{4}{C} \punkt{0}{4}{D} \punkt{0}{-3}{S} \end{ksys}
\teil Eine Pyramide hat den quadratischen Boden $ABCD$ mit der Seite $2$ cm; die Spitze $S$ liegt $1{,}5$ cm senkrecht über $A$. Im Gitter sind der Boden und die Spitze des Seitendreiecks $ABS$ markiert. Zeichne die drei übrigen Seitendreiecke ein. \hfill (Abitur 2024 GK)

\begin{ksys}[xmin=-2.5,xmax=5.5,ymin=-2.5,ymax=5.5] \punkt{0}{0}{A} \punkt{2}{0}{B} \punkt{2}{2}{C} \punkt{0}{2}{D} \punkt{0}{-1.5}{S} \end{ksys}
\teil Die Pyramide $ABCS$ hat die Grundfläche $A(4 | 1 | 2)$, $B(1 | 5 | 2)$, $C(-1 | -1 | 2)$ und die Spitze $S(4 | -1 | 6)$. Eingezeichnet sind die Projektionen $A'$ und $B'$ in die $x_1x_2$-Ebene. Ergänze $C'$ und $S'$ und entscheide: Liegt der Höhenfußpunkt innerhalb der Grundfläche? \hfill (Abitur 2018 GK)

\begin{ksys3} \rpunkt*{4}{1}{0}{A'} \rpunkt*{1}{5}{0}{B'} \end{ksys3}
\teil Die Pyramide $ABCS$ hat die Grundfläche $A(5 | -1 | 1)$, $B(1 | 5 | 1)$, $C(-1 | -2 | 1)$ und die Spitze $S(1 | 1 | 6)$. Eingezeichnet sind die Projektionen $A'$ und $B'$ in die $x_1x_2$-Ebene. Ergänze $C'$ und $S'$ und entscheide: Liegt der Höhenfußpunkt innerhalb der Grundfläche? \hfill (Abitur 2018 GK)

\begin{ksys3} \rpunkt*{5}{-1}{0}{A'} \rpunkt*{1}{5}{0}{B'} \end{ksys3}
\end{teile}
\end{aufgabe}

% Anlauf (ohne Prüfkennung)
\begin{aufgabe}{Streckenlänge, Mittelpunkt, Teilpunkte – Anlauf}
\begin{teile}
\teil Wo muss der Haken an der Stange zwischen $A$ und $B$ hängen, damit er genau in der Mitte ist? Was ist gefragt? \\ \kreuz{eine Länge} \\ \kreuz{ein Punkt} \\ \kreuz{eine Figur}
\teil $A(2 | 1 | 3)$, $B(4 | 4 | 9)$: Länge von $AB$? \leerfeld LE
\teil Ein Zeltmast endet in $S(2 | 2 | 6)$; vier gleich lange Leinen führen zu Heringen am Boden, einer davon steckt in $H(4 | 5 | 0)$ (1 LE = 1 m). Je Leine kommen $40$ cm für die Knoten dazu. Wie viel Leine braucht man insgesamt? \leerfeld[m]
\end{teile}
\end{aufgabe}

\begin{aufgabe}{Streckenlänge, Mittelpunkt, Teilpunkte – Prüfungsaufgaben}
\begin{teile}
\teil Ein Gummiseil ist von $A(3 | 1 | 2)$ nach $B(9 | 9 | 2)$ gespannt (1 LE = 1 m) und dabei um $25\,\%$ länger als ungespannt. Wie lang ist es ungespannt? \hfill (Abitur 2020 GK) \\ \leerfeld[m]
\teil Ein Seil hängt belastet von $A(1 | 1 | 10)$ bis $B(3 | 7 | 1)$ (1 LE = 1 m) und ist dabei $10\,\%$ länger als unbelastet. Unbelastete Länge? \hfill (Abitur 2020 GK) \\ \leerfeld[m]
\teil Eine Feder zwischen $A(2 | 3 | 1)$ und $B(3 | 7 | 9)$ (1 LE = 1 cm) ist $50\,\%$ länger als in Ruhe. Länge in Ruhe? \hfill (Abitur 2020 GK) \\ \leerfeld[cm]
\teil Eine Pyramide ist symmetrisch zur $x_1x_3$-Ebene; $D(1 | 3 | 0)$ ist eine Grundecke, $S(3 | 0 | 6)$ die Spitze. Zeige: $W(2 | 1{,}5 | 3)$ ist der Mittelpunkt der Kante $DS$. Gib den Mittelpunkt $T$ der zu $DS$ symmetrischen Kante an. \hfill (Abitur 2026 GK) \\ T( \leerfeld | \leerfeld | \leerfeld )
\teil Eine Dachgaube ist symmetrisch zur $x_2x_3$-Ebene; $D(6 | 1 | 0)$ ist eine Ecke, $S(0 | 5 | 8)$ die Firstspitze. Zeige: $W(3 | 3 | 4)$ ist der Mittelpunkt der Kante $DS$. Gib den Mittelpunkt $T$ der zu $DS$ symmetrischen Kante an. \hfill (Abitur 2026 GK) \\ T( \leerfeld | \leerfeld | \leerfeld )
\teil Ein Kristall ist eine Doppelpyramide, symmetrisch zur $x_1x_2$-Ebene; $D(7 | 4 | 0)$ ist eine Ecke, $S(1 | 2 | 4)$ die obere Spitze. Zeige: $W(4 | 3 | 2)$ ist der Mittelpunkt der Kante $DS$. Gib den Mittelpunkt $T$ der zu $DS$ symmetrischen Kante an. \hfill (Abitur 2026 GK) \\ T( \leerfeld | \leerfeld | \leerfeld )
\teil Eine Strebe führt von $P(9 | 0 | 3)$ nach $Q(0 | 3 | 0)$; zwei Schellen teilen sie in drei gleich lange Stücke. Koordinaten der Schellen? \hfill (Abitur 2024 GK)
\teil Ein Kletterseil ist von Haken $A(1 | 4 | 2)$ zu Haken $C(11 | -1 | 7)$ gespannt. Der Karabiner $G$ teilt es so, dass $|AG| : |GC| = 2 : 3$. Koordinaten von $G$? \hfill (Abitur 2023 GK) \\ G( \leerfeld | \leerfeld | \leerfeld )
\teil Eine Stange reicht von $A(2 | 0 | 8)$ bis $C(6 | 8 | 0)$. Eine Leuchte $G$ sitzt so, dass $|AG| : |GC| = 3 : 1$. Koordinaten von $G$? \hfill (Abitur 2023 GK) \\ G( \leerfeld | \leerfeld | \leerfeld )
\teil $A(3 | 5 | 2)$ und $B(5 | 6 | 4)$ liegen auf einer Geraden $g$. Gib einen Punkt $D \ne B$ auf $g$ an, der von $A$ so weit entfernt ist wie $B$. \hfill (Abitur 2020 GK) \\ D( \leerfeld | \leerfeld | \leerfeld )
\teil $A(1 | -2 | 4)$ und $B(4 | 2 | 4)$ liegen auf einer Geraden $g$. Gib einen Punkt $D \ne B$ auf $g$ an, der von $A$ so weit entfernt ist wie $B$. \hfill (Abitur 2020 GK) \\ D( \leerfeld | \leerfeld | \leerfeld )
\teil $A(0 | 3 | 5)$ und $B(2 | 1 | 6)$ liegen auf einer Geraden $g$. Gib einen Punkt $D \ne B$ auf $g$ an, der von $A$ so weit entfernt ist wie $B$. \hfill (Abitur 2020 GK) \\ D( \leerfeld | \leerfeld | \leerfeld )
\end{teile}
\end{aufgabe}

\begin{aufgabe}{Streckenlänge, Mittelpunkt, Teilpunkte – Prüfungsaufgaben – weiter}
\begin{teile}
\teil $A(2 | -1 | 2)$; der Punkt $C$ ist festgelegt durch $\overrightarrow{OC} = 3 \cdot \overrightarrow{OA}$. Länge der Strecke $AC$? \hfill (Abitur 2018 GK) \\ \leerfeld
\teil $A(2 | 3 | 6)$; der Punkt $C$ ist festgelegt durch $\overrightarrow{OC} = - \overrightarrow{OA}$. Länge der Strecke $AC$? \hfill (Abitur 2018 GK) \\ \leerfeld
\teil $A(4 | 4 | 2)$; der Punkt $C$ ist festgelegt durch $\overrightarrow{OC} = \frac{1}{2} \cdot \overrightarrow{OA}$. Länge der Strecke $AC$? \hfill (Abitur 2018 GK) \\ \leerfeld
\teil Ein Tunnel verläuft geradlinig von $A(10 | 30 | 3)$ nach $B(40 | 70 | 3)$ (1 LE = 10 m). Ein Zug fährt ab $A$ mit $60$ km/h, gleichzeitig ein zweiter ab $B$ mit $90$ km/h. Welchen Weg im Tunnel legt der Zug ab $B$ zurück, bis sich beide begegnen? \hfill (Abitur 2019 GK) \\ \leerfeld[m]
\teil Ein Radweg führt geradlinig von $P(2 | 1 | 0)$ nach $Q(8 | 9 | 0)$ (1 LE = 100 m). Anna startet in $P$ mit $15$ km/h, Ben gleichzeitig in $Q$ mit $10$ km/h. Wie weit ist Anna gefahren, wenn sie sich treffen, und wo ist das? \hfill (Abitur 2019 GK)
\end{teile}
\end{aufgabe}

% Anlauf (ohne Prüfkennung)
\begin{aufgabe}{Dreiecke nachweisen – Anlauf}
\begin{teile}
\teil Parallelogramm $ABCD$: Welcher Vektor ist gleich $\overrightarrow{AB}$? \\ \kreuz{$\overrightarrow{CD}$} \\ \kreuz{$\overrightarrow{DC}$} \\ \kreuz{$\overrightarrow{BC}$}
\end{teile}
\end{aufgabe}

\begin{aufgabe}{Dreiecke nachweisen – Prüfungsaufgaben}
\begin{teile}
\teil Zeige: Das Dreieck $A(2 | 0 | 1)$, $B(6 | 2 | 1)$, $C(4 | 1 | 5)$ ist gleichschenklig. Nenne Basis und Schenkel. \hfill (Abitur 2023 GK)
\teil Zeige: Das Dreieck $A(0 | 2 | 3)$, $B(4 | 2 | 1)$, $C(2 | 5 | 2)$ ist gleichschenklig. Nenne Basis und Schenkel. \hfill (Abitur 2023 GK)
\teil Zeige: Das Dreieck $P(1 | 4 | 0)$, $Q(5 | 2 | 0)$, $R(5 | 7 | 2)$ ist gleichschenklig. Nenne Basis und Schenkel. \hfill (Abitur 2023 GK)
\teil Zeige: Das Dreieck $A(2 | 1 | 4)$, $B(4 | 5 | 2)$, $C(5 | 3 | 5)$ ist gleichschenklig. Nenne Basis und Schenkel. \hfill (Abitur 2023 GK)
\teil Zeige: Das Dreieck $A(0 | 3 | 1)$, $B(4 | 1 | 5)$, $C(4 | 6 | 3)$ ist gleichschenklig. Nenne Basis und Schenkel. \hfill (Abitur 2023 GK)
\teil Prüfe, ob das Dreieck $A(4 | 2 | 2)$, $B(2 | 4 | 2)$, $C(2 | 2 | 4)$ gleichseitig ist. \hfill (Abitur 2023 GK)
\teil Prüfe, ob das Dreieck $A(5 | 2 | 2)$, $B(2 | 5 | 2)$, $C(2 | 2 | 6)$ gleichseitig ist. \hfill (Abitur 2020 GK)
\teil Prüfe, ob das Dreieck $A(7 | 3 | 1)$, $B(3 | 7 | 1)$, $C(3 | 3 | 5)$ gleichseitig ist. \hfill (Abitur 2023 GK)
\teil $A(1 | 2 | 0)$, $B(5 | 6 | 0)$ und $C(3 | 4 | p)$ mit $p > 0$. Zeige: Das Dreieck $ABC$ ist für jedes $p$ gleichschenklig mit der Basis $AB$. \hfill (Abitur 2022 GK)
\teil $A(3 | 0 | 3)$, $B(3 | 6 | 1)$ und $C(p | 3 | 2)$. Zeige: Das Dreieck $ABC$ ist für jedes $p \ne 3$ gleichschenklig mit der Basis $AB$. \hfill (Abitur 2022 GK)
\teil $A(-1 | 2 | 2)$, $B(3 | 2 | -2)$ und $C(1 | p | 0)$. Zeige: Das Dreieck $ABC$ ist für jedes $p$ gleichschenklig. Welche Seite ist die Basis? \hfill (Abitur 2022 GK)
\teil Ein Quader hat die Ecken $A(9 | 0 | 0)$, $B(9 | 9 | 0)$ und $C(0 | 9 | 0)$ in der $x_1x_2$-Ebene. Begründe: Das Dreieck $ABC$ ist rechtwinklig und gleichschenklig. Flächeninhalt? \hfill (Abitur 2018 GK)
\end{teile}
\end{aufgabe}

\begin{aufgabe}{Dreiecke nachweisen – Prüfungsaufgaben – weiter}
\begin{teile}
\teil Grundfläche $A(6 | 1 | 0)$, $B(1 | 4 | 0)$, $C(1 | -2 | 0)$; Deckkante $E(1 | 3 | 2)$, $F(1 | -1 | 2)$. Zeige: $ABC$ ist gleichschenklig. Begründe: $EF$ ist parallel zur Grundfläche. \hfill (Abitur 2023 GK)
\teil Grundfläche $A(6 | 4 | 3)$, $B(2 | 1 | 3)$, $C(2 | 7 | 3)$; Deckkante $G(1 | 3 | 5)$, $H(1 | 5 | 5)$. Zeige: $ABC$ ist gleichschenklig. Begründe: $GH$ ist parallel zur Grundfläche. \hfill (Abitur 2023 GK)
\teil Grundfläche $A(5 | 6 | 0)$, $B(3 | 1 | 0)$, $C(7 | 1 | 0)$; Deckkante $P(4 | 2 | 3)$, $Q(6 | 2 | 3)$. Zeige: $ABC$ ist gleichschenklig. Begründe: $PQ$ ist parallel zur Grundfläche. \hfill (Abitur 2023 GK)
\teil $D(8 | 0 | 2)$, $E(0 | 6 | 2)$, $F(0 | 0 | 2)$ und $M(4 | 3 | 2)$. Begründe: $D$, $E$ und $F$ liegen auf einem Kreis mit dem Mittelpunkt $M$. \hfill (Abitur 2024 GK)
\teil Begründe: $D(2 | 2 | 5)$, $E(6 | 6 | 1)$ und $F(6 | 2 | 5)$ liegen auf einem Kreis, dessen Mittelpunkt die Mitte von $DE$ ist. \hfill (Abitur 2024 GK)
\teil Begründe: $P(-1 | 3 | 0)$, $Q(5 | 3 | 0)$ und $R(2 | 3 | 3)$ liegen auf einem Kreis, dessen Mittelpunkt die Mitte von $PQ$ ist. \hfill (Abitur 2024 GK)
\teil In der $x_1x_2$-Ebene liegen $A(0 | 0)$, $B(6 | 0)$ und $C(3 | 9)$ auf einem Kreis. Koordinaten seines Mittelpunkts $M$? \hfill (Abitur 2019 GK) \\ M( \leerfeld | \leerfeld )
\teil In der $x_1x_2$-Ebene liegen $A(-2 | 0)$, $B(6 | 0)$ und $C(2 | 8)$ auf einem Kreis. Koordinaten seines Mittelpunkts $M$? \hfill (Abitur 2019 GK) \\ M( \leerfeld | \leerfeld )
\teil In der $x_1x_2$-Ebene liegen $A(1 | 1)$, $B(7 | 1)$ und $C(4 | 10)$ auf einem Kreis. Koordinaten seines Mittelpunkts $M$? \hfill (Abitur 2019 GK) \\ M( \leerfeld | \leerfeld )
\end{teile}
\end{aufgabe}

\begin{aufgabe}{Dreiecke nachweisen – Prüfungsaufgaben – weiter}
\begin{teile}
\teil Dreieck $A(2 | 1 | 2)$, $B(4 | 5 | 4)$, $C(5 | 3 | 1)$ mit $|CA| = |CB|$. Gesucht ist ein Punkt $D \ne C$ mit $|DA| = |DB|$ und $|DA| = |CA|$. Koordinaten von $D$? \hfill (Abitur 2023 GK) \\ D( \leerfeld | \leerfeld | \leerfeld )
\teil Dreieck $A(0 | 2 | 3)$, $B(4 | 2 | -1)$, $C(4 | 3 | 3)$ mit $|CA| = |CB|$. Gesucht ist ein Punkt $D \ne C$ mit $|DA| = |DB|$ und $|DA| = |CA|$. Koordinaten von $D$? \hfill (Abitur 2023 GK) \\ D( \leerfeld | \leerfeld | \leerfeld )
\teil $A(3 | 1 | 2)$ und $M(1 | 3 | 1)$. Gib zwei Punkte $B$ und $C$ an, sodass $A$, $B$ und $C$ denselben Abstand zu $M$ haben und ein gleichschenkliges Dreieck bilden. \hfill (Abitur 2024 GK)
\teil $A(2 | 0 | 5)$ und $M(4 | 2 | 4)$. Gib zwei Punkte $B$ und $C$ an, sodass $A$, $B$ und $C$ denselben Abstand zu $M$ haben und ein gleichschenkliges Dreieck bilden. \hfill (Abitur 2024 GK)
\end{teile}
\end{aufgabe}

% Anlauf (ohne Prüfkennung)
\begin{aufgabe}{Vierecke nachweisen – Anlauf}
\begin{teile}
\teil Nachzuweisen: $ABCD$ ist ein Rechteck. Was gehört zum Steckbrief? \\ \kreuz{vier gleich lange Seiten} \\ \kreuz{Parallelogramm und ein rechter Winkel} \\ \kreuz{ein Paar paralleler Seiten}
\end{teile}
\end{aufgabe}

\begin{aufgabe}{Vierecke nachweisen – Prüfungsaufgaben}
\begin{teile}
\teil Zeige: $A(2 | 1 | 0)$, $B(4 | 3 | 1)$, $C(2 | 4 | 3)$, $D(0 | 2 | 2)$ bilden ein Rechteck $ABCD$. \hfill (Abitur 2023 GK)
\teil $A(1 | 5 | p)$, $B(4 | 5 | p)$, $C(4 | 1 | 5)$ und $D(1 | 1 | 5)$ mit einer natürlichen Zahl $p$. Zeige: $ABCD$ ist für jedes $p$ ein Rechteck. Länge von $BC$ in Abhängigkeit von $p$? \hfill (Abitur 2025 GK)
\teil Zeige: $A(0 | 1 | 3)$, $B(2 | 3 | 4)$, $C(3 | 5 | 6)$, $D(1 | 3 | 5)$ bilden eine Raute $ABCD$. \hfill (Abitur 2022 GK)
\teil Eine Dachfläche hat die Ecken $A(0 | 1 | 1)$, $B(4 | 1 | 4)$, $C(4 | 6 | 4)$, $D(0 | 6 | 1)$. Zeige: $ABCD$ ist eine Raute. \hfill (Abitur 2022 GK)
\teil Zeige: $A(1 | 2 | 2)$, $B(2 | 6 | 10)$, $C(6 | 7 | 18)$, $D(5 | 3 | 10)$ bilden eine Raute $ABCD$. \hfill (Abitur 2022 GK)
\teil Zeige: $A(2 | 2 | 2)$, $B(5 | 3 | 2)$, $C(6 | 5 | 4)$, $D(3 | 4 | 4)$ bilden ein Parallelogramm, aber kein Rechteck. \hfill (Abitur 2022 GK)
\teil Zeige: $A(1 | 3 | 2)$, $B(3 | 4 | 4)$, $C(5 | 2 | 5)$, $D(3 | 1 | 3)$ bilden eine Raute, aber kein Quadrat. \hfill (Abitur 2020 GK)
\teil Zeige: $A(1 | 0 | 1)$, $B(1 | 3 | 5)$, $C(4 | 7 | 5)$, $D(4 | 4 | 1)$ bilden eine Raute, aber kein Quadrat. \hfill (Abitur 2020 GK)
\teil Eine Kletterwand hat die Ecken $A(4 | 1 | 3)$, $B(1 | 4 | 3)$, $F(1 | 7 | 0)$ und $E(7 | 1 | 0)$ (1 LE = 1 m). Zeige: $ABFE$ ist ein Trapez, in dem zwei gegenüberliegende Seiten gleich lang sind. \hfill (Abitur 2018 GK)
\teil Ein Würfel mit Kante $8$ hat eine Ecke im Ursprung und liegt im ersten Oktanten. $P(8 | 0 | 2)$, $Q(3 | 8 | 0)$, $R(0 | 8 | 3)$ und $S(2 | 0 | 8)$ liegen auf seinen Kanten. Begründe: $PQRS$ ist ein Trapez mit zwei gleich langen Seiten, und $PS$ ist doppelt so lang wie $QR$. \hfill (Abitur 2019 GK)
\end{teile}
\end{aufgabe}

\begin{aufgabe}{Vierecke nachweisen – Prüfungsaufgaben – weiter}
\begin{teile}
\teil Eine Terrasse ist das Viereck $A(1 | -3 | 0)$, $B(7 | -1 | 0)$, $C(7 | 1 | 0)$, $D(1 | 3 | 0)$ (1 LE = 1 m). Zeige: $ABCD$ ist ein Trapez. Flächeninhalt? \hfill (Abitur 2026 GK) \\ \leerfeld[m²]
\teil Ein Beet ist das Viereck $E(-2 | 1 | 0)$, $F(6 | 1 | 0)$, $G(3 | 5 | 0)$, $H(1 | 5 | 0)$ (1 LE = 1 m). Zeige: $EFGH$ ist ein Trapez. Flächeninhalt? \hfill (Abitur 2026 GK) \\ \leerfeld[m²]
\teil Eine Glasplatte ist das Viereck $K(0 | -3 | 1)$, $L(4 | -1 | 1)$, $M(4 | 1 | 1)$, $N(0 | 3 | 1)$ (1 LE = 1 dm). Zeige: $KLMN$ ist ein Trapez. Flächeninhalt? \hfill (Abitur 2026 GK) \\ \leerfeld dm²
\teil Zeige: Das Viereck $K(0 | 0 | 2)$, $L(2 | 3 | 2)$, $M(6 | 0 | 2)$, $N(2 | -3 | 2)$ ist ein Drachenviereck. \hfill (Abitur 2023 GK)
\teil Zeige: Das Viereck $K(1 | 2 | 0)$, $L(1 | 5 | 2)$, $M(1 | 2 | 8)$, $N(1 | -1 | 2)$ ist ein Drachenviereck. \hfill (Abitur 2023 GK)
\teil Zeige: Das Viereck $P(6 | 1 | 1)$, $Q(6 | 5 | 1)$, $R(6 | 4 | 4)$, $S(6 | 1 | 5)$ ist ein Drachenviereck. \hfill (Abitur 2023 GK)
\teil Eine ebene Rasenfläche hat die Ecken $A(2 | 1 | 0)$, $B(9 | 1 | 0)$, $C(8 | 4 | 0{,}5)$, $D(8 | 7 | 1)$ und $E(2 | 7 | 1)$ (1 LE = 1 m). Zeige: $AE$ ist parallel zu $CD$, und $CD$ und $DE$ bilden einen rechten Winkel. \hfill (Abitur 2021 GK)
\teil Zeige: Im Viereck $A(1 | -1 | 4)$, $B(3 | 1 | 5)$, $C(1 | 2 | 2{,}5)$, $D(0 | 1 | 2)$ ist $AB$ parallel zu $DC$, und bei $A$ liegt ein rechter Winkel. \hfill (Abitur 2021 GK)
\teil Zeige: Im Viereck $A(0 | 2 | 1)$, $B(6 | 2 | 9)$, $C(7 | 3 | 2)$, $D(4 | 3 | -2)$ ist $AB$ parallel zu $DC$, und bei $A$ liegt ein rechter Winkel. \hfill (Abitur 2021 GK)
\teil $P(2 | 1 | 0)$ und $Q(2 | 1 | 5)$ sind benachbarte Ecken eines Quadrats in der Ebene $3x_1 + 4x_2 = 10$. Gib eine weitere Ecke an. \hfill (Abitur 2022 GK)
\teil $P(3 | 1 | 4)$ und $Q(3 | 6 | 4)$ sind benachbarte Ecken eines Quadrats in der Ebene $4x_1 - 3x_3 = 0$. Gib eine weitere Ecke an. \hfill (Abitur 2022 GK)
\end{teile}
\end{aufgabe}

\begin{aufgabe}{Vierecke nachweisen – Prüfungsaufgaben – weiter}
\begin{teile}
\teil Das Viereck $ABCD$ hat die Ecken $A(0 | 0 | 0)$, $B(3 | 1 | 0)$, $C(0 | 7 | 0)$ und $D(-2 | 3 | 0)$. $B$ wird parallel zur $x_2$-Achse nach $B'$ verschoben, sodass $AB'CD$ bei $B'$ einen rechten Innenwinkel hat; $b$ ist die gesuchte Koordinate. Erläutere den Ansatz $(3 | b | 0) \circ (3 | b - 7 | 0) = 0$. \hfill (Abitur 2025 GK)
\teil Das Viereck $ABCD$ hat die Ecken $A(0 | 0 | 0)$, $B(2 | 4 | 0)$, $C(9 | 0 | 0)$ und $D(4 | -3 | 0)$. $B$ wird parallel zur $x_1$-Achse nach $B'$ verschoben, sodass $AB'CD$ bei $B'$ einen rechten Innenwinkel hat; $b$ ist die gesuchte Koordinate. Erläutere den Ansatz $(b | 4 | 0) \circ (b - 9 | 4 | 0) = 0$. \hfill (Abitur 2025 GK)
\teil Die Raute $A(5 | 3 | 1)$, $B(2 | 3 | 5)$, $C(-1 | 3 | 1)$, $D(2 | 3 | -3)$ hat einen Inkreis, der alle vier Seiten berührt. Begründe: $P(3{,}92 | 3 | 2{,}44)$ ist einer der Berührpunkte. \hfill (Abitur 2020 GK)
\teil Die Raute $A(1 | 1 | 6)$, $B(4 | 1 | 2)$, $C(1 | 1 | -2)$, $D(-2 | 1 | 2)$ hat einen Inkreis, der alle vier Seiten berührt. Begründe: $P(2{,}92 | 1 | 3{,}44)$ ist einer der Berührpunkte. \hfill (Abitur 2020 GK)
\teil Eine Wand hat die Form des Parallelogramms $A(2 | 1 | 0)$, $B(9 | 1 | 0)$, $C(10 | 1 | 4)$, $D(3 | 1 | 4)$ (1 LE = 1 m). Ein Farbstreifen von $A$ nach $P(6 | 1 | 4)$ auf der Seite $DC$ teilt sie. Welcher Teil ist größer? Begründe ohne Flächenberechnung. \hfill (Abitur 2022 GK)
\teil Eine rechteckige Wand hat die Ecken $A(1 | 3 | 0)$, $B(1 | 9 | 0)$, $C(1 | 9 | 5)$, $D(1 | 3 | 5)$ (1 LE = 1 m). Ein Streifen von $B$ nach $Q(1 | 3 | 2)$ auf der Seite $AD$ teilt sie. Welcher Teil ist größer? Begründe ohne Flächenberechnung. \hfill (Abitur 2022 GK)
\end{teile}
\end{aufgabe}

% Anlauf (ohne Prüfkennung)
\begin{aufgabe}{Körper und Drehungen – Anlauf}
\begin{teile}
\teil Ein Quader steht auf der $x_1x_2$-Ebene; seine Deckenecke $G$ hat $x_3 = 4$. Was sagt das über die Deckfläche? \\ \kreuz{Sie liegt in der $x_1x_2$-Ebene.} \\ \kreuz{Sie liegt parallel zur $x_1x_2$-Ebene in der Höhe $4$.} \\ \kreuz{Sie steht senkrecht auf der $x_1x_2$-Ebene.}
\end{teile}
\end{aufgabe}

\begin{aufgabe}{Körper und Drehungen – Prüfungsaufgaben}
\begin{teile}
\teil Gerades Prisma $ABCDEF$ ($D$ über $A$, $E$ über $B$, $F$ über $C$) mit $A(2 | 1 | 0)$, $B(7 | 1 | 0)$, $C(2 | 4 | 0)$, $D(2 | 1 | 3)$: $F$? \hfill (Abitur 2026 GK) \\ F( \leerfeld | \leerfeld | \leerfeld )
\teil Quader $ABCDEFGH$ (Boden $ABCD$, $E$ über $A$, $F$ über $B$, $G$ über $C$, $H$ über $D$) mit $A(2 | 1 | 0)$, $B(6 | 1 | 0)$, $E(2 | 1 | 5)$, $H(2 | 4 | 5)$: $G$? \hfill (Abitur 2024 GK) \\ G( \leerfeld | \leerfeld | \leerfeld )
\teil Quader $ABCDEFGH$ (Boden $ABCD$, $E$ über $A$, $F$ über $B$, $G$ über $C$, $H$ über $D$) mit $A(0 | 2 | 1)$, $B(4 | 2 | 1)$, $E(0 | 2 | 4)$, $H(0 | 7 | 4)$: $G$? \hfill (Abitur 2024 GK) \\ G( \leerfeld | \leerfeld | \leerfeld )
\teil Gerades Prisma $ABCDEF$ ($D$ über $A$, $E$ über $B$, $F$ über $C$) mit $A(4 | 1 | 2)$, $B(4 | 6 | 2)$, $C(8 | 1 | 2)$, $D(4 | 1 | 7)$: $F$? \hfill (Abitur 2026 GK) \\ F( \leerfeld | \leerfeld | \leerfeld )
\teil Eine Pyramide hat als Grundfläche ein rechtwinkliges Dreieck mit den Katheten $6$ cm und $8$ cm; die Spitze liegt $5$ cm senkrecht über der Ecke mit dem rechten Winkel. Wähle geeignete Koordinaten der vier Ecken (1 LE = 1 cm). \hfill (Abitur 2020 GK)
\teil Eine Halle ist ein Quader mit quadratischem Grundriss (Seite $15$ m) und $7$ m Höhe. $A$ liegt im Ursprung, $B$ auf der $x_1$-Achse, $D$ auf der $x_2$-Achse, die Deckenecken $E$, $F$, $G$, $H$ über $A$, $B$, $C$, $D$ (1 LE = 1 m). Koordinaten von $G$ und $H$? \hfill (Abitur 2020 GK)
\teil Ein Würfel mit Kante $2$ hat eine Ecke im Ursprung, alle Ecken haben Koordinaten $\ge 0$. Er wird um $(-1 | -1 | 1)$ verschoben. Welche Ecke hat danach $x_1 < 0$, $x_2 > 0$ und $x_3 > 2$? \hfill (Abitur 2024 GK) \\ ( \leerfeld | \leerfeld | \leerfeld )
\teil Ein Quader mit den Kanten $4$, $2$ und $3$ entlang der $x_1$-, $x_2$- und $x_3$-Achse hat eine Ecke im Ursprung, alle Ecken haben Koordinaten $\ge 0$. Er wird um $(-2 | 1 | -1)$ verschoben. Welche Ecke hat danach $x_1 > 0$, $x_2 < 2$ und $x_3 < 0$? \hfill (Abitur 2024 GK) \\ ( \leerfeld | \leerfeld | \leerfeld )
\teil Ein Quader mit den Kanten $3$, $5$ und $2$ entlang der $x_1$-, $x_2$- und $x_3$-Achse hat eine Ecke im Ursprung, alle Ecken haben Koordinaten $\ge 0$. Er wird um $(-2 | -3 | -1)$ verschoben. Welche Ecke hat danach $x_1 > 0$, $x_2 < 0$ und $x_3 > 0$? \hfill (Abitur 2024 GK) \\ ( \leerfeld | \leerfeld | \leerfeld )
\teil Eine Pyramide hat die Grundfläche $A(4 | -1 | 2)$, $B(1 | 3 | 2)$, $C(-1 | 0 | 2)$ und die Spitze $S(2 | 1 | 7)$. Begründe: Die Grundfläche ist parallel zur $x_1x_2$-Ebene. Höhe der Pyramide? \hfill (Abitur 2018 GK) \\ h = \leerfeld
\teil Eine Pyramide hat die Grundfläche $A(2 | 2 | -1)$, $B(5 | 0 | -1)$, $C(3 | 5 | -1)$ und die Spitze $S(3 | 2 | 3)$. Begründe: Die Grundfläche ist parallel zur $x_1x_2$-Ebene. Höhe der Pyramide? \hfill (Abitur 2018 GK) \\ h = \leerfeld
\teil Eine Pyramide hat die Grundfläche $A(1 | 0 | 3)$, $B(4 | 2 | 3)$, $C(0 | 5 | 3)$ und die Spitze $S(2 | 2 | 9)$. Begründe: Die Grundfläche ist parallel zur $x_1x_2$-Ebene. Höhe der Pyramide? \hfill (Abitur 2018 GK) \\ h = \leerfeld
\end{teile}
\end{aufgabe}

\begin{aufgabe}{Körper und Drehungen – Prüfungsaufgaben – weiter}
\begin{teile}
\teil Ein Dachgeschoss hat einen rechteckigen Boden von $10$ m mal $8$ m. Die Dachflächen steigen von den beiden Traufwänden (Höhe $0{,}5$ m) gleichmäßig zum First über der Mitte (Höhe $5{,}5$ m). Flächen mit weniger als $2$ m Raumhöhe zählen nicht zur Wohnfläche. Welcher Anteil des Bodens zählt nicht? \hfill (Abitur 2019 GK) \\ \leerfeld \%
\teil Ein Dachgeschoss hat einen rechteckigen Boden von $14$ m mal $10$ m. Die Dachflächen steigen von den beiden Traufwänden (Höhe $1$ m) gleichmäßig zum First über der Mitte (Höhe $5$ m). Flächen mit weniger als $1{,}8$ m Raumhöhe zählen nicht zur Wohnfläche. Welcher Anteil des Bodens zählt nicht? \hfill (Abitur 2019 GK) \\ \leerfeld \%
\teil Ein Dachgeschoss hat einen rechteckigen Boden von $15$ m mal $10$ m. Die Dachflächen steigen von den beiden Traufwänden (Höhe $0{,}4$ m) gleichmäßig zum First über der Mitte (Höhe $2{,}9$ m). Flächen mit weniger als $1$ m Raumhöhe zählen nicht zur Wohnfläche. Welcher Anteil des Bodens zählt nicht? \hfill (Abitur 2019 GK) \\ \leerfeld \%
\teil Das Dreieck $A(5 | 3 | 0)$, $B(3 | 5 | 0)$, $C(3 | 3 | 4)$ wird um die Kante $AB$ gedreht, bis $C$ in der $x_1x_2$-Ebene liegt, mit $x_1 < 4$. Koordinaten des gedrehten Punktes $C'$? \hfill (Abitur 2023 GK) \\ C'( \leerfeld | \leerfeld | \leerfeld )
\teil Das Dreieck $A(2 | 1 | 0)$, $B(2 | 9 | 0)$, $C(5 | 5 | 4)$ wird um die Kante $AB$ gedreht, bis $C$ in der $x_1x_2$-Ebene liegt, mit $x_1 > 2$. Koordinaten des gedrehten Punktes $C'$? \hfill (Abitur 2023 GK) \\ C'( \leerfeld | \leerfeld | \leerfeld )
\teil Eine Pyramide hat den quadratischen Boden $ABCD$ mit der Seite $4$ in der $x_1x_2$-Ebene, $A$ im Ursprung, $B$ auf der $x_1$-Achse, $D$ auf der $x_2$-Achse; die Spitze $S$ liegt $6$ senkrecht über $A$. Ein Würfel hat drei Seitenflächen in den Koordinatenebenen und eine Ecke auf der Kante $CS$. Kantenlänge des Würfels? Begründe, dass kein Punkt des Würfels außerhalb der Pyramide liegt. \hfill (Abitur 2024 GK)
\teil Eine Pyramide hat den quadratischen Boden $ABCD$ mit der Seite $3$ in der $x_1x_2$-Ebene, $A$ im Ursprung, $B$ auf der $x_1$-Achse, $D$ auf der $x_2$-Achse; die Spitze $S$ liegt $6$ senkrecht über $A$. Ein Würfel hat drei Seitenflächen in den Koordinatenebenen und eine Ecke auf der Kante $CS$. Kantenlänge des Würfels? Begründe, dass kein Punkt des Würfels außerhalb der Pyramide liegt. \hfill (Abitur 2024 GK)
\teil Ein Körper hat den rechteckigen Boden $A(2 | 1 | 0)$, $B(6 | 1 | 0)$, $C(6 | 7 | 0)$, $D(2 | 7 | 0)$ und die Firstkante $P(4 | 2 | 3)$, $Q(4 | 5 | 3)$; die Dreiecke $ABP$ und $CDQ$ und die Vierecke $BCQP$ und $DAPQ$ begrenzen ihn. Die Ebenen $x_2 = s$ mit $1 < s < 7$ schneiden ihn. Wie viele Ecken hat die Schnittfläche? Für welche $s$ haben die Schnittflächen dieselbe Form und denselben Flächeninhalt? \hfill (Abitur 2024 GK)
\teil Ein Körper hat den rechteckigen Boden $A(0 | 1 | 0)$, $B(6 | 1 | 0)$, $C(6 | 9 | 0)$, $D(0 | 9 | 0)$ und die Firstkante $P(3 | 3 | 4)$, $Q(3 | 6 | 4)$; die Dreiecke $ABP$ und $CDQ$ und die Vierecke $BCQP$ und $DAPQ$ begrenzen ihn. Die Ebenen $x_2 = s$ mit $1 < s < 9$ schneiden ihn. Wie viele Ecken hat die Schnittfläche? Für welche $s$ haben die Schnittflächen dieselbe Form und denselben Flächeninhalt? \hfill (Abitur 2024 GK)
\teil Ein Würfel $ABCDEFGH$ mit Kante $4$ hat die Ecke $A$ im Ursprung, $B$ auf der $x_1$-Achse, $D$ auf der $x_2$-Achse; $E$ liegt über $A$, $G$ über $C$. Die Ebene durch $E$, $G$ und $S_t(t | 0 | 0)$ mit $0 \le t \le 4$ schneidet ihn. Eckenzahl der Schnittfläche in Abhängigkeit von $t$? Für welche $t$ hat sie genau zwei Symmetrieachsen? \hfill (Abitur 2024 GK)
\teil Ein Quader $ABCDEFGH$ mit quadratischem Boden (Seite $6$) und Höhe $3$ hat die Ecke $A$ im Ursprung, $B$ auf der $x_1$-Achse, $D$ auf der $x_2$-Achse; $E$ liegt über $A$, $G$ über $C$. Die Ebene durch $E$, $G$ und $S_t(t | 0 | 0)$ mit $0 \le t \le 6$ schneidet ihn. Eckenzahl der Schnittfläche in Abhängigkeit von $t$? Für welche $t$ hat sie genau zwei Symmetrieachsen? \hfill (Abitur 2024 GK)
\end{teile}
\end{aufgabe}
